\documentclass[twoside,letterpaper,12pt]{report}
\usepackage[spanish,backgroundcolor=green,textsize=small]{todonotes}
\usepackage[margin=2cm]{geometry}
\usepackage[utf8x]{inputenc}
\usepackage[spanish]{babel}
\usepackage{plain}
\usepackage{amsmath}
\usepackage{graphicx}
\usepackage{wrapfig}

\usepackage{parskip}
\setlength{\parindent}{15pt}

\usepackage{hyperref}
\hypersetup{
  colorlinks=true,
  linkcolor=black,          % color of internal links (change box color with linkbordercolor)
  citecolor=black,        % color of links to bibliography
  filecolor=black,      % color of file links
  urlcolor=black           % color of external links
}

\title{
Proyecto de Grado\\[0.5cm]
Instalación, Implementación y Uso de las librerías boost y su aplicación en recorrido de grafos para la supercomputadora Apolo}
\vspace{1cm}

\author{
	Autor:\\[0.3cm]
	Sergio Andrés Monsalve Castañeda\\
	Código: 200410061010\\
	smonsal3@eafit.edu.co\\
	semonsalve@gmail.com\\[0.7cm]
	Asesor: \\[0.3cm]
	Juan Guillermo Lalinde Pulido\\
	jlalinde@eafit.edu.co\\
	2619500 ext 9588%\\[1cm]
}

\date{
	\today \\[0.5cm]
	\includegraphics[width=0.25\textwidth]{aux/logo_eafit}
}


\begin{document}
\maketitle

\tableofcontents


\begin{abstract}

El crecimiento de la capacidad computacional abre nuevos campos de aplicación. La teoría de grafos es una de las herramientas más importantes que ofrece las ciencias de la computación para la resolución de diferentes tipos de problemas. Aunque los algoritmos para grafos tienen una complejidad alta en general, existen algoritmos paralelos. En este proyecto se pretende instalar en Apolo, el supercomputador de EAFIT, la infraestructura básica para que un programador, utilizando las librerías BOOST para C++, pueda desarrollar rápidamente aplicaciones que aprovechen su capacidad computacional. 

El crecimiento de la capacidad computacional abre nuevos campos de aplicación. 

La teoría de grafos es una de las herramientas más importantes que ofrece las ciencias de la computación para la resolución de diferentes tipos de problemas. 

Aunque los algoritmos para grafos tienen una complejidad alta en general, existen algoritmos paralelos. En este proyecto se pretende instalar en Apolo, el supercomputador de EAFIT, la infraestructura básica para que un programador, utilizando las librerías BOOST para C++, pueda desarrollar rápidamente aplicaciones que aprovechen su capacidad computacional. 


Comunidad Científica
Comunidad académica
Estudiantes de Ingeniería de Sistemas
Estudiantes de Ingenieria Matematica


\end{abstract}

\newpage

\renewcommand{\abstractname}{Agradecimientos}
\begin{abstract}
Juan Guillermo Lalinde, Juan David Pineda, John Jairo Silva, Alejandro Gómez, Mateo Gómez, Jaime Pérez, John Mario Gutiérrez.
\end{abstract}

\newpage

\chapter{Introducción}

\section{Objetivos}

\begin{itemize}
	\item Construir un manual para la instalación de las librerías de boost
	\item Generar Documentación sobre el manejo de las librerías de boost
	\item Crear códigos de ejemplo para la ejecución de las librerías de boost	
\end{itemize}



\section{Alcance y productos}

Descripción de las Librerías de Boost
Librerías de Utilización de Grafos Instaladas en Apolo
Manual de Instalacion de librerias Para Boost
Ejemplo en código documentado y manual de como hacer un programa utilizando boost para grafos.
Introducción al manejo de Grafos con boost (representación conceptual)

\section{Usabilidad}

Utilizar los GPS de los taxis para hacer análisis de la malla vial de la ciudad. 
Movilización y Rutas óptimas para el recorrido de Bomberos, Ambulancias, Policía para la atención de emergencias y desastres


Movilidad
Seguridad
Educación 
Tecnología

\chapter{Marco de Referencia}
\label{ChapRef}

\section{Apolo} Super Computadores
\section{Grafos}
\section{MPI}
\section{Boost}

\begin{quotation}
The Parallel Boost Graph Library is an extension to the Boost Graph Library (BGL) for parallel and distributed computing. It offers distributed graphs and graph algorithms to exploit coarse-grained parallelism along with parallel algorithms that exploit fine-grained parallelism, while retaining the same interfaces as the (sequential) BGL. Code written using the sequential BGL should be easy to parallelize with the parallel BGL. Visitors new to the Parallel BGL should read our architectural overview.\cite{wwwboost} 
\end{quotation} 


\section{Metaprogramación}
\subsection{Metaprogramación en C++}


\chapter{Instalación}



\section{Virtualización de un cluster}[imagen de apolo]

\section{MPI}

\subsection{OpenMPI}
\subsection{Mpich}

\section{Boost Graph Library}

	build mpi
	build GraphParallel
	build serilization

Boost and Eclipse
	configuración
	
\subsection{Parallel Boost Graph Library (Parallel BGL) }
\subsubsection{Seriallization}
\subsubsection{Regex}
\subsubsection{Parallel}
\subsubsection{MPI}



\subsection{Instalacion de Boost en Apolo}


\subsection{makefile}

Cómo ejecutar el script de installación 

que se debe modificar en el makefile\cite{Wall2000}

instalar rpm de boost:
Cuando se formatea o reinicia máquina
Cuando se instala normalmente.



\chapter{Implementación}

Ejemplo Programa Corriendo en Apolo

\chapter{Resultados}

\section{Sin paralelizar}
\section{10 Nodos}
\section{100 Nodos}



\chapter{Glosario}
\label{chapGlosario}

\begin{description}
	\item[HPC:] (High Performance Computing) Computación de alto desempeño.
	\item[Data Center:]
\end{description}

\newpage

\bibliography{Tesis}

\bibliographystyle{plain}	
aca
\cite{czarnecki2000generative}
\cite{wwwboost}
\cite{stroustrup2013c++}
\cite{andrei2001modern}
\cite{Wall2000}
\cite{Boost}
\cite{Karniadakis}
\cite{Kernighan1988}


\todo[inline,caption={TODO}]{
}

\newpage

\appendix
\chapter{Repositorio del código}

Para una copia del código utilizado dirigirse a: \url{https://github.com/smonsalve/tesis.git}

OVA de la maquina virtual con Rocks instalado y dos nodos.

Input data para Grafo con X Nodos. y Y vertices.

Makefile de Instalación.

\end{document}

